\section{Graph Theory}
\subsection{Graph Condense (Kosaraju SCC)}
\begin{minted}{c++}
struct SCC {
  int n, comps;
  vector<vector<int>> adj, rev_adj, scc_adj;
  vector<int> order, comp;
  vector<bool> vis;
  SCC(int _n) : n(_n), comps(0), adj(_n), rev_adj(_n),
    comp(_n, -1), vis(_n, false) {}
  void add_edge(int u, int v) {
    adj[u].push_back(v);
    rev_adj[v].push_back(u);
  }
  void dfs1(int u) {
    vis[u] = true;
    for (int v : adj[u]) if (!vis[v]) dfs1(v);
    order.push_back(u);
  }
  void dfs2(int u, int id) {
    comp[u] = id;
    for (int v : rev_adj[u])
      if (comp[v] == -1) dfs2(v, id);
  }
  void build() {
    for (int i = 0; i < n; i++) if (!vis[i]) dfs1(i);
    reverse(order.begin(), order.end());
    for (int u : order) {
      if (comp[u] == -1) dfs2(u, comps++);
    }
    scc_adj.resize(comps);
    for (int u = 0; u < n; u++) {
      for (int v : adj[u]) {
        if (comp[u] != comp[v])
          scc_adj[comp[u]].push_back(comp[v]);
      }
    }
  }
};
\end{minted}
\subsection{Articulation Points \texorpdfstring{\&}{&} Bridges}
\begin{minted}{c++}
struct Articulation {
  int n, timer;
  vector<vector<int>> adj;
  vector<int> tin, low;
  vector<bool> is_art;
  vector<pair<int, int>> bridges;
  Articulation(int _n) : n(_n), timer(0), adj(_n),
    tin(_n, -1), low(_n, -1), is_art(_n, false) {}
  void add_edge(int u, int v) {
    adj[u].push_back(v);
    adj[v].push_back(u);
  }
  void dfs(int u, int p = -1) {
    tin[u] = low[u] = timer++;
    int children = 0;
    for (int v : adj[u]) {
      if (v == p) continue;
      if (tin[v] != -1) {
        low[u] = min(low[u], tin[v]);
      } else {
        children++;
        dfs(v, u);
        low[u] = min(low[u], low[v]);
        if (low[v] >= tin[u] && p != -1)
          is_art[u] = true;
        if (low[v] > tin[u])
          bridges.push_back({min(u,v), max(u,v)});
      }
    }
    if (p == -1 && children > 1) is_art[u] = true;
  }
  void build() {
    for (int i = 0; i < n; i++)
      if (tin[i] == -1) dfs(i);
  }
};
\end{minted}
\subsection{Hopcroft-Karp (Bipartite Matching)}
\begin{minted}{c++}
struct HopcroftKarp {
  int n, m;
  vector<vector<int>> g;
  vector<int> dist, pairU, pairV;
  static const int INF = 1e9;
  HopcroftKarp(int n_ = 0, int m_ = 0) {
    n = n_; m = m_;
    g.assign(n, {});
    pairU.assign(n, -1);
    pairV.assign(m, -1);
    dist.assign(n, 0);
  }
  void add_edge(int u, int v) { g[u].push_back(v); }
  bool bfs() {
    queue<int> q;
    for (int u = 0; u < n; ++u) {
      if (pairU[u] == -1) {
        dist[u] = 0; q.push(u);
      } else dist[u] = INF;
    }
    bool foundAug = false;
    while (!q.empty()) {
      int u = q.front(); q.pop();
      for (int v : g[u]) {
        int u2 = pairV[v];
        if (u2 == -1) foundAug = true;
        else if (dist[u2] == INF) {
          dist[u2] = dist[u] + 1; q.push(u2);
        }
      }
    }
    return foundAug;
  }
  bool dfs(int u) {
    for (int v : g[u]) {
      int u2 = pairV[v];
      if (u2 == -1 || (dist[u2] == dist[u] + 1
      && dfs(u2))) {
        pairU[u] = v; pairV[v] = u; return true;
      }
    }
    dist[u] = INF; return false;
  }
  int max_matching() {
    int matching = 0;
    while (bfs()) {
      for (int u = 0; u < n; ++u)
        if (pairU[u] == -1 && dfs(u)) ++matching;
    }
    return matching;
  }
  vector<pair<int,int>> matching_edges() const {
    vector<pair<int,int>> res;
    for (int v = 0; v < m; ++v)
      if (pairV[v] != -1)
        res.emplace_back(pairV[v], v);
    return res;
  }
};
\end{minted}
\subsection{Dinic's Algorithm (Max Flow)}
\begin{minted}{c++}
const int64_t inf = 1LL << 61;
struct Dinic {
  struct edge { int to, rev;
    int64_t flow, w; };
  int n, s, t;
  vector<int> d, done;
  vector<vector<edge>> g;
  Dinic() {}
  Dinic(int _n) { n = _n + 10, g.resize(n); }
  void add_edge(int u, int v, int64_t w) {
    edge a = {v, (int)g[v].size(), 0, w};
    edge b = {u, (int)g[u].size(), 0, 0};
    g[u].push_back(a), g[v].push_back(b);
  }
  bool bfs() {
    d.assign(n, -1), d[s] = 0;
    queue<int> Q; Q.push(s);
    while (Q.size()) {
      int u = Q.front(); Q.pop();
      for (auto &e : g[u]) {
        int v = e.to;
        if (!~d[v] && e.flow < e.w)
          d[v] = d[u] + 1, Q.push(v);
      }
    }
    return d[t] != -1;
  }
  int64_t dfs(int u, int64_t flow) {
    if (u == t) return flow;
    for (int &i = done[u]; i < (int)g[u].size(); i++) {
      edge &e = g[u][i];
      if (e.w <= e.flow) continue;
      int v = e.to;
      if (d[v] == d[u] + 1) {
        int64_t nw = dfs(v, min(flow, e.w - e.flow));
        if (nw > 0) {
          e.flow += nw;
          g[v][e.rev].flow -= nw;
          return nw;
        }
      }
    }
    return 0;
  }
  int64_t max_flow(int _s, int _t) {
    s = _s, t = _t;
    int64_t flow = 0;
    while (bfs()) {
      done.assign(n, 0);
      while (int64_t nw = dfs(s, inf)) flow += nw;
    }
    return flow;
  }
};
\end{minted}
\subsection{Biconnected Comp. \texorpdfstring{\&}{&} Bridge Tree}
\begin{minted}{c++}
struct BCC {
  int n, timer, comps;
  vector<vector<int>> adj, tree;
  vector<int> tin, low, comp;
  vector<bool> is_bridge, vis;
  vector<pair<int, int>> edges;
  BCC(int _n) : n(_n), timer(0), comps(0), adj(_n),
    tin(_n, -1), low(_n, -1), comp(_n, -1),
    vis(_n, false) {}
  void add_edge(int u, int v) {
    adj[u].push_back(edges.size());
    adj[v].push_back(edges.size());
    edges.push_back({u, v});
    is_bridge.push_back(false);
  }
  void dfs(int u, int p = -1) {
    vis[u] = true;
    tin[u] = low[u] = timer++;
    for (int id : adj[u]) {
      int v = edges[id].first ^
              edges[id].second ^ u;
      if (v == p) continue;
      if (vis[v]) low[u] = min(low[u], tin[v]);
      else {
        dfs(v, u);
        low[u] = min(low[u], low[v]);
        if (low[v] > tin[u]) is_bridge[id] = true;
      }
    }
  }
  void build_tree() {
    for (int i = 0; i < n; i++) if (!vis[i]) dfs(i);
    vector<int> q;
    for (int i = 0; i < n; i++) {
      if (comp[i] != -1) continue;
      q.push_back(i); comp[i] = comps++;
      while (!q.empty()) {
        int u = q.back(); q.pop_back();
        for (int id : adj[u]) {
          int v = edges[id].first ^
                  edges[id].second ^ u;
          if (comp[v] == -1 && !is_bridge[id]) {
            comp[v] = comp[u]; q.push_back(v);
          }
        }
      }
    }
    tree.resize(comps);
    for (int i = 0; i < (int)edges.size(); i++) {
      if (is_bridge[i]) {
        int u = comp[edges[i].first],
            v = comp[edges[i].second];
        tree[u].push_back(v);
        tree[v].push_back(u);
      }
    }
  }
};
\end{minted}
\subsection{Binary Lifting (LCA \texorpdfstring{\&}{&} K-th Ancestor)}
\begin{minted}{c++}
struct LCA {
  int n, LOG;
  vector<vector<int>> adj, up;
  vector<int> depth;
  LCA(int _n) : n(_n), LOG(__lg(_n) + 2), adj(_n),
    up(_n, vector<int>(LOG, -1)), depth(_n, 0) {}
  void add_edge(int u, int v) {
    adj[u].push_back(v);
    adj[v].push_back(u);
  }
  void dfs(int u, int p = -1, int d = 0) {
    up[u][0] = p;
    depth[u] = d;
    for (int i = 1; i < LOG; i++) {
      if (up[u][i - 1] != -1)
        up[u][i] = up[up[u][i - 1]][i - 1];
    }
    for (int v : adj[u])
      if (v != p) dfs(v, u, d + 1);
  }
  int lift(int u, int k) {
    for (int i = 0; i < LOG; i++) {
      if ((k >> i) & 1) u = up[u][i];
      if (u == -1) break;
    }
    return u;
  }
  int query(int u, int v) {
    if (depth[u] < depth[v]) swap(u, v);
    u = lift(u, depth[u] - depth[v]);
    if (u == v) return u;
    for (int i = LOG - 1; i >= 0; i--) {
      if (up[u][i] != up[v][i]) {
        u = up[u][i]; v = up[v][i];
      }
    }
    return up[u][0];
  }
};
\end{minted}
\subsection{Tree Flattening (Euler Tour)}
\begin{minted}{c++}
struct EulerTour {
  int n, timer;
  vector<vector<int>> adj;
  vector<int> in, out, flat;
  EulerTour(int _n) : n(_n), timer(0), adj(_n + 1),
    in(_n + 1), out(_n + 1) {}
  void add_edge(int u, int v) {
    adj[u].push_back(v);
    adj[v].push_back(u);
  }
  void dfs(int u, int p = 0) {
    in[u] = ++timer;
    flat.push_back(u);
    for (int v : adj[u])
      if (v != p) dfs(v, u);
    out[u] = timer;
  }
  void build(int root = 1) { dfs(root); }
};
\end{minted}
\subsection{Heavy Light Decomposition (HLD)}
\begin{minted}{c++}
struct HLD {
  int n, cur_pos = 1;
  vector<int> arr, parent, size, depth, heavy,
  head, pos;
  vector<vector<int>> adj;
  BIT<ll> segtree;
  HLD(int n) : n(n), arr(n + 1), parent(n + 1),
  size(n + 2), depth(n + 2), heavy(n + 2, 0),
  head(n + 2), pos(n + 2), adj(n + 2),
  segtree(BIT<ll>(n + 2)) {}
  void add_edge(int u, int v) {
    adj[u].push_back(v), adj[v].push_back(u);
  }
  void dfs(int u, int p) {
    parent[u] = p; size[u] = 1;
    int max_sz = 0;
    for (int v : adj[u]) {
      if (v == p) continue;
      depth[v] = depth[u] + 1;
      dfs(v, u);
      size[u] += size[v];
      if (size[v] > max_sz)
        max_sz = size[v], heavy[u] = v;
    }
  }
  void decompose(int u, int h) {
    head[u] = h; pos[u] = cur_pos++;
    segtree.add(pos[u], arr[u]);
    if (heavy[u]) decompose(heavy[u], h);
    for (int v : adj[u])
      if (v != parent[u] && v != heavy[u])
        decompose(v, v);
  }
  int query(int u, int v) {
    int res = 0;
    while (head[u] != head[v]) {
      if (depth[head[u]] < depth[head[v]])
        swap(u, v);
      res += segtree.sum(pos[head[u]], pos[u]);
      u = parent[head[u]];
    }
    if (depth[u] > depth[v]) swap(u, v);
    res += segtree.sum(pos[u], pos[v]);
    return res;
  }
  void addVal(int idx, int delta) {
    segtree.add(pos[idx], delta);
    arr[idx] += delta;
  }
  void editVal(int idx, int value) {
    segtree.add(pos[idx], value - arr[idx]);
    arr[idx] = value;
  }
};
\end{minted}
\subsection{Centroid Decomposition}
\begin{minted}{c++}
vector<vector<int>> adj;
vector<bool> is_removed;
vector<int> subtree_size;
int get_subtree_size(int node, int parent = -1) {
  subtree_size[node] = 1;
  for (int child : adj[node]) {
    if (child == parent || is_removed[child])
      continue;
    subtree_size[node] +=
    get_subtree_size(child, node);
  }
  return subtree_size[node];
}
int get_centroid(int node, int tree_size,
int parent = -1) {
  for (int child : adj[node]) {
    if (child == parent || is_removed[child])
      continue;
    if (subtree_size[child] * 2 > tree_size)
      return get_centroid(child, tree_size, node);
  }
  return node;
}
void build_centroid_decomp(int node = 0) {
  int centroid = get_centroid(node,
  get_subtree_size(node));
  // do something with centroid
  is_removed[centroid] = true;
  for (int child : adj[centroid]) {
    if (is_removed[child]) continue;
    build_centroid_decomp(child);
  }
}
\end{minted}
\subsection{Centroid Decomposition (Ahsan's)}
\begin{minted}{c++}
int tot = 0;
vector<int> sz(n, 1), removed(n, 0), cenpar(n, -1);
auto dfs_sz = [&](auto &&self, int node, int pr)
-> void {
  tot += 1; sz[node] = 1;
  for (auto &son : g[node]) {
    if (son == pr || removed[son]) continue;
    self(self, son, node);
    sz[node] += sz[son];
  }
};
auto find_cen = [&](auto &&self, int node, int pr)
-> int {
  for (auto &son : g[node]) {
    if (son == pr || removed[son]) continue;
    if (sz[son] > tot / 2)
      return self(self, son, node);
  }
  return node;
};
auto decompose = [&](auto &&self, int node, int pr)
-> void {
  tot = 0; dfs_sz(dfs_sz, node, pr);
  int cen = find_cen(find_cen, node, pr);
  removed[cen] = 1; cenpar[cen] = pr;
  for (auto &son : g[cen]) {
    if (son == pr || removed[son]) continue;
    self(self, son, cen);
  }
};
decompose(decompose, 0, -1);
\end{minted}
\subsection{LCA and Distance}
\begin{minted}{c++}
vector<vector<int>> t, parents;
vector<int> level;
int mx;
void dfs(int node, int parent, int cLevel) {
  level[node] = cLevel;
  parents[0][node] = parent;
  for (auto &x : t[node])
    if (x != parent) dfs(x, node, cLevel + 1);
}
int kthParent(int a, int k) {
  for (int i = 0; i < mx; i++) {
    if (a == -1) return a;
    if (k & (1 << i)) a = parents[i][a];
  }
  return a;
}
int lca(int a, int b) {
  if (level[a] > level[b])
    a = kthParent(a, level[a] - level[b]);
  else b = kthParent(b, level[b] - level[a]);
  if (a == b) return a;
  for (int i = mx - 1; i >= 0; i--) {
    if (parents[i][a] != parents[i][b]) {
      a = parents[i][a]; b = parents[i][b];
    }
  }
  return parents[0][a];
}
int distance(int a, int b) {
  return level[a] + level[b] - 2 * level[lca(a, b)];
}
\end{minted}
\subsection{DSU on Tree (Sack)}
$O(N \log N)$ offline subtree queries.
\begin{minted}{c++}
int timer = 0;
vector<int> sz, in, out, flat, big;
void dfs_sz(int u, int p) {
  sz[u] = 1, big[u] = 0, in[u] = ++timer,
  flat[timer] = u;
  for (int v : adj[u]) {
    if (v != p) {
      dfs_sz(v, u);
      sz[u] += sz[v];
      if (sz[v] > sz[big[u]]) big[u] = v;
    }
  }
  out[u] = timer;
}
void add(int u, int val) { /* update DS */ }
void dfs_sack(int u, int p, bool keep) {
  for (int v : adj[u])
    if (v != p && v != big[u]) dfs_sack(v, u, 0);
  if (big[u]) dfs_sack(big[u], u, 1);
  for (int v : adj[u])
    if (v != p && v != big[u])
      for (int i = in[v]; i <= out[v]; i++)
        add(flat[i], 1);
  add(u, 1);
  // answer queries for subtree u here
  if (!keep)
    for (int i = in[u]; i <= out[u]; i++)
      add(flat[i], -1);
}
// Init: sz/in/out/flat/big.assign(n+1, 0);
// dfs_sz(1, 0); dfs_sack(1, 0, 1);
\end{minted}
